%% ma: L10-B12-0065
%% lop: 10
%% chuong: 5
%% bai: 12
%% dang: 10.12.8
%% loai: DS
%% muc: [NB, TH, TH, TH]
%% dap_an: "ĐSĐĐ"
%% nguon: "chuyen-de-cac-so-dac-trung-cua-mau-so-lieu-khong-ghep-nhom-toan-10-chuong-trinh-moi.pdf (GV. Phan Nhật Linh) (P01), Bài 01 Phần II Câu 9"
%% kien_thuc: "Lớp 10 Bài 12 (số gần đúng, sai số tuyệt đối, độ chính xác, sai số tương đối, quy tròn số)"
%% trang_thai: chinh-thuc
%% ghi_chu: "Tự duyệt theo yêu cầu giáo viên 22/10/2026."
\begin{ex}
Đường kính của một đồng hồ cát là $8{,}52$ m. Cho giá trị $\pi=3{,}141592654\ldots$ Xét tính đúng sai của các khẳng định sau:
\choiceTF{\True Giá trị gần đúng của $\pi$ làm tròn đến hàng phần nghìn là $3{,}142$}{Chu vi đồng hồ cát bằng đúng $26{,}8$ m}{\True Bán kính của đồng hồ cát làm tròn đến hàng phần mười là $4{,}3$ m}{\True Dùng giá trị gần đúng của $\pi$ làm tròn đến hàng phần trăm để tính chu vi, rồi làm tròn kết quả đến hàng phần mười, được $26{,}8$ m}
\loigiai{a) Đúng.

b) Sai: chu vi $C=8{,}52\pi\approx26{,}766$ m, không bằng $26{,}8$ m.

c) Đúng: $r=4{,}26$ m $\approx4{,}3$ m.

d) Đúng: $3{,}14\cdot8{,}52=26{,}7528\approx26{,}8$ m.}
\end{ex}
